\documentclass[polish]{article}
\usepackage{graphicx} % Required for inserting images

\usepackage[polish]{babel}
\usepackage{polski}
\usepackage{placeins}
\usepackage{dirtytalk}
% Set page size and margins
% Replace `letterpaper' with `a4paper' for UK/EU standard size
\usepackage{csquotes}
% Useful packages
\usepackage{float}
\usepackage{amsmath}
\usepackage{graphicx}
\usepackage{hyperref}
\usepackage{graphicx}
\usepackage{listings}
\usepackage{xcolor}
\usepackage{amssymb}
\usepackage{float}
\definecolor{codegreen}{rgb}{0,0.6,0}
\definecolor{codegray}{rgb}{0.5,0.5,0.5}
\definecolor{codepurple}{rgb}{0.58,0,0.82}
\definecolor{backcolour}{rgb}{0.95,0.95,0.92}

\lstdefinestyle{mystyle}{
    backgroundcolor=\color{backcolour},   
    commentstyle=\color{codegreen},
    keywordstyle=\color{magenta},
    numberstyle=\tiny\color{codegray},
    stringstyle=\color{codepurple},
    basicstyle=\ttfamily\footnotesize,
    breakatwhitespace=false,         
    breaklines=true,                 
    captionpos=b,                    
    keepspaces=true,                 
    numbers=left,                    
    numbersep=5pt,                  
    showspaces=false,                
    showstringspaces=false,
    showtabs=false,                  
    tabsize=2
}
\lstset{style=mystyle}
\hypersetup{
    colorlinks=true,
    linktoc=all     %set to all if you want both sections and subsections linked
}
\usepackage{mdframed}
\title{Zagadnienia do 1. kolokwium}
\author{}
\date{}

\begin{document}

\maketitle

\section{Teoria}
\begin{enumerate}
    \item Przestrzeń zdarzeń elementarnych, $\sigma$-algebra zdarzeń. 
    \item Miara probabilistyczna (definicja i własności).
    \item Prawdopodobieństwo warunkowe (definicja i własności).
    \item Niezależność zdarzeń i niezależność warunkowa (w obu przypadkach definicja i własności).
    \item Układ zupełny, twierdzenie o prawdopodobieństwie całkowitym (z dowodem).
    \item Twierdzenie Bayesa (z dowodem).
    \item Zmienna losowa (definicja, rozkład prawdopodobieństwa -- twierdzenie).
    \item Dystrybuanta (definicja, warunki i własności).
    \item Gęstość zmiennej losowej (definicja, warunki i własności).
    \item Wartość oczekiwana, momenty (w tym momenty centralne), wariancja, odchylenie standardowe, kwantyle.
    \item Funkcje generujące momenty.
    \item Metoda dystrybuanty i metoda transformacji dla zmiennych jednowymiarowych.
    \item Rozkłady: dwumianowy, geometryczny, Poissona, jednostajny, wykładniczy, normalny.
\end{enumerate}

\end{document}

